\reportchapter{源码、文档与贡献}

\section{重构后源码}

重构后源码仓库链接：

\begin{quote}
\url{https://github.com/dyx131313/SmartSpar_BikeHub}
\end{quote}

主要目录：

\begin{itemize}
  \item 后端入口：\code{backend/run.py}
  \item 后端应用：\code{backend/app}
  \item 后端服务层：\code{backend/app/services}
  \item 后端工具：\code{backend/app/utils}
  \item 前端入口：\code{frontend/bikehub/src}
  \item 前端页面：\code{frontend/bikehub/src/features}
\end{itemize}

\section{文档清单}

本次交付文档包括：

\begin{itemize}
  \item \code{docs/LOCAL_RUNBOOK.md}：本地运行手册。
  \item \code{docs/REFACTORING_RECORD.md}：重构过程记录。
  \item \code{docs/TEST_REPORT.md}：测试报告。
  \item \code{docs/REFACTORING_REPORT_DRAFT.md}：Markdown 版报告草稿。
  \item \code{docs/refactoring_report_tex/}：基于同济 TeX 模板的正式报告工程。
  \item \code{docs/refactoring_report_tex/figures/evidence/}：页面截图、架构图、覆盖矩阵和命令行验证图。
  \item \code{docs/refactoring_report_tex/evidence/}：lint、build、pytest、接口烟测的原始输出。
\end{itemize}

\section{原设计文档与使用说明修订}

作业要求中的“文档：修订原设计文档和使用说明书”已在本次交付中单独落实。原后端文档原本分散在 \code{backend/docs/} 下，部分接口说明与当前服务层、权限控制和预测模块不一致；重构后统一整理到根目录 \code{docs/}，并按提交材料用途拆分为设计说明、接口说明、运行手册和预测指南。

\begin{longtable}{p{0.24\textwidth}p{0.34\textwidth}p{0.28\textwidth}}
\toprule
文档类型 & 修订后文件 & 修订内容 \\
\midrule
原设计文档 & \code{docs/SmartSpar BikeHub 后端系统完整文档.md} & 补充分层架构、服务层职责、权限装饰器、统一响应、预测注册表和调度服务说明。 \\
接口设计说明 & \code{docs/智慧共享单车调度系统 - 后端API文档.md} & 更新登录、用户、站点、任务、需求预测、系统时间和反馈接口的当前路径与返回结构。 \\
使用说明书 & \code{docs/LOCAL_RUNBOOK.md} & 补充 Windows 本地启动、MySQL 3307 初始化、测试账号、前后端端口和常见故障处理。 \\
预测使用说明 & \code{docs/需求预测运行指南.md} & 说明模型文件位置、future 数据读取、训练依赖降级和预测接口验收方式。 \\
\bottomrule
\end{longtable}

因此，最终提交不仅包含重构报告，也包含修订后的设计文档和使用说明书；第 5 章测试报告和附录命令清单用于支撑这些文档的可复现性。

\section{贡献说明}

小组名称为“葫芦娃除BUG”。成员分工和贡献度已单独整理在封面后的“小组贡献说明”页，作为报告第二页展示。

\section{总结}

本次重构将 SmartSpar BikeHub 从一个可以演示但工程基线不稳定的项目，整理为可以本地启动、可以构建、可以测试、可以说明的提交版本。后端通过服务层、权限装饰器、统一响应和路径配置降低了维护复杂度；前端通过构建修复、运营中台化设计和关键路由修复提升了可用性；测试和文档则保证重构结果可复现、可验收。
